\pdfoutput=1
\documentclass[12pt]{article}
\usepackage{amsmath, amssymb, amsthm}
\usepackage{graphicx}
\usepackage[colorlinks=true, linkcolor=blue, citecolor=blue, urlcolor=blue]{hyperref}
\usepackage[numbers,sort&compress]{natbib}
\usepackage{geometry}
\geometry{margin=1in}

\newtheorem{theorem}{Theorem}[section]
\newtheorem{lemma}[theorem]{Lemma}
\newtheorem{proposition}[theorem]{Proposition}
\newtheorem{corollary}[theorem]{Corollary}
\theoremstyle{definition}
\newtheorem{definition}[theorem]{Definition}
\theoremstyle{remark}
\newtheorem{remark}[theorem]{Remark}

\title{Game-Theoretic Framing of the Triadic Balance Condition}
\author{Jeffrey W. Milton\\\\[4pt]\small Clarity Coalition}
\date{2026}

\begin{document}
\maketitle

\begin{abstract}
The tholonic triad assigns three functionally distinct roles (Negotiation $N$,
Definition $D$, and Contribution $C$) to a three-variable recurrence whose branches
converge to classical constants: $\pi/4$, $\varphi$, $e$, $\sqrt{2}$, and $\ln 2$.
We recast this system as a two-player alternating-move game in which $D$ and $C$ act
as strategic agents whose opposing objectives (constraining versus accumulating) generate
a dynamical trajectory for the payoff state $N$.  The triadic balance condition is
identified as the equilibrium reached when the marginal contributions of $D$ and $C$ to
$N$ are equal and opposite, a condition that reduces to the fixed-point equation of the
underlying recurrence.  We prove that this equilibrium is a pure-strategy Nash
equilibrium of the associated stage game, that the diagonal-invariance and swap-symmetry
theorems of the original framework admit direct characterizations in terms of zero-sum
and symmetric subgame structure, and that the three traversal classes (Advancing,
Self-redefined, Fixed) correspond to distinct information structures in the game:
exogenous shocks, endogenous state evolution, and constant-strategy play, respectively.
The framework provides a new lens: five classical constants emerge not merely as limits
of recurrences but as equilibrium payoffs of a minimal strategic interaction.
\end{abstract}

\tableofcontents

